\section{Synthesis of the second question: duplication, and the \texorpdfstring{$b$}{b}-field between manifolds}
\label{s6:sec:synthesis}

We return to the author's two remarks of \cref{s0:sec:question}: the $b$-particle is a deal breaker when two manifolds or more are in
question, and some process makes a unitary energy result in a duplication, from $\Cl(3,3)$ to $\Cl(6,6)$. The verdicts are collected in
\cref{s6:tab:verdict}. As in \cref{s5:sec:synthesis}, a rate becomes a rest energy only through the added identification $mc^2=\hbar\kappa$.

\subsection{How a unitary energy results in a duplication}

\paragraph{The mathematical process is forced on the algebra and chosen on the energy.} Doubling the vector space forces the graded tensor
product $\Cl(V\oplus V)=\Cl(V)\hat\otimes\Cl(V)$. Its Klein factor is determined up to conjugation, it is the Jordan--Wigner sign, and the
spinor module of the duplicate is the Fock space of six fermionic modes (\cref{d5:sec:algebra}; inherited). The energy is not forced. The
undivided energy $(\C^3,E_6,U(3))$ embeds isometrically in $(\C^3\otimes\C^2,E_6\oplus E_6,U(6))$. The Lorentz-covariant embeddings form the
three-sphere of a lossless beam splitter, and those that also respect $\eta$ form one circle, the copy rotation (\cref{d5:thm:dup}). Read
literally as a multiplication of the energy, the phrase fails: a linear duplication divides the undivided energy between the copies, and a
quadratic copy $z\mapsto z\otimes z$ is excluded by the no-cloning theorem (inherited).

\paragraph{The physical processes} that take one copy to two are listed in \cref{d5:tab:processes}. Two of them concern mass. A duplication that
preserves the Lorentz action cannot turn a single null field into a massive pair (\cref{d6:prop:covariant}). Making a pair needs either a
second, non-collinear null field, which gives the kinematics of $\gamma\gamma\to e^+e^-$ (\cref{d6:thm:pair,d6:ex:bw}), or a non-covariant third
body that takes the momentum, as a nucleus does in pair production (\cref{d6:ex:R1}). All of this is inherited kinematics in the arena's
variables. The closure never creates a pair (\cref{d6:cor:nopair}).

\paragraph{What the duplicate can do that one copy cannot} (fixed, partly restated from \cite{ES17}). It carries an $\eta$-isometric complex
structure and a compact, Lorentz-invariant copy charge (\cref{d5:prop:cx}), the copy algebra $\so{2,\C}$ on two binding frames
(\cref{d5:prop:so6C}), a non-rigid Lorentz-invariant vacuum on $S^4$ and the imaginary unit $J=J_1\otimes\omega$ (\cref{d5:prop:vacuum}). In
the closure, massive constituents (non-null momentum fields, $E_0(\pi^{(k)})\neq0$) can turn at an exactly constant rate on a codimension-two balance surface (\cref{d6:thm:balance}). This last
point is the precise dynamical sense in which the duplicate does what one copy cannot, and it needs no $b$-field. The duplicate also loses two
things: spin one-half under the diagonal interior action (two fermions make a boson, \cref{d5:prop:spin}), and the uniqueness of the photon
direction.

\paragraph{The caveat.} The rest energy of each copy separately is not invariant under the copy rotation, so it is gauge dependent if that
rotation is gauged. With equal inertias, the closure's own symmetry maps
every balanced pair to one null constituent (\cref{d6:prop:gauge}). A statement that a mass has been created in one copy therefore needs a
structure that distinguishes the copies. The symmetric closure and the gauged copy rotation provide none. Unequal rates do provide one, and
that is where the second remark enters.

\subsection{The \texorpdfstring{$b$}{b}-field between manifolds}

\paragraph{It exists, and where it lives depends on the Lorentz action} (fixed; \cref{d7:prop:exist}). If the Lorentz algebra acts on all $n$
binding frames in the same way, the binding-compatible fields between them form $\so{n,\C}$, with compact part the copy rotations $\so n$:
abelian for two manifolds, non-abelian ($\so3$) from three on. If the Lorentz algebra acts on one manifold only, that manifold is decoupled
from every binding-compatible field, and the fields act among the internal manifolds, as the paper's $\so4_\Delta$ does.

\paragraph{The remark is true for transmission} (fixed). Between manifolds with unequal rates a compact, binding-compatible inter-manifold field cannot be removed
along a worldline. It acts as a Hermitian coupling in the Schr\"odinger form of the frame flow. It shifts and spreads the lines, confines a frame without
rate of its own by a seesaw, and at special ratios creates a resonance with secular growth (\cref{d5:thm:bfield2,d7:thm:schr}). From three
manifolds on it carries an invariant that two manifolds lack: the sign of $B_{12}B_{23}B_{31}$, odd under interior time reversal and fixing the
sense in which energy circulates. This is the arena's analogue of Jarlskog's invariant (\cref{d7:thm:three}). On one manifold a field that commutes with
the binding root is removable, and none of this happens. A field that does not commute with it can there be non-removable and shift the lines
(\cref{d4:thm:screen}), but it confines no frame without rate, and a connection alone binds nothing, on one manifold or on several
(\cref{d4:prop:alone,d5:rem:dealbreaker}).

\paragraph{The remark is false for creation} (fixed). The rates are the spectrum of the self-adjoint part of the shift, which transforms as a
tensor, so no field, between manifolds or on one, creates or changes a rate. The field redistributes the lines of the motion, the frequencies, at fixed sum. With equal
rates every binding-compatible field is removable, and that includes the non-abelian $\so4_\Delta$. With zero rates the field produces only drift
(\cref{d7:thm:schr}). The deal breaker therefore presupposes a differentiation of the rates. The paper's one-rotor closure does not supply it,
since it imposes a common rate. Under a common Lorentz action the closure admits a rotor acting by a non-central element of $\uu n$, whose
copy-field part and rate part are independent; such a closure admits unequal rates together with a copy field (\cref{d7:prop:rotor}; constructed,
not selected).

\paragraph{Isospin one-half does not cross between two or three manifolds that share the Lorentz algebra} (fixed). There the copy isospin is
integral on the whole tower module, because a binding-compatible copy rotation turns an even number of planes. The Yang--Mills doublet acting on two copies of the
undivided energy is a connection in one of its three directions only. Its other two directions are self-adjoint, of the same type as the
binding (\cref{d7:prop:isospin}). On three manifolds doublets need the spectator action. There, unequal internal rates break $\so4_\Delta=\su2\oplus\su2$ to a
diagonal $\so3$ and make three directions physical (\cref{d7:prop:so4unequal}). From four manifolds on the copy rotations themselves contain
$\so4\cong\su2\oplus\su2$, and the copy index carries doublets without a spectator (\cref{d7:rem:verdict}(d)).

\subsection{\texorpdfstring{Back to $E=mc^2$}{Back to E = mc2}}

The second half sharpens the verdict of \cref{s5:sec:synthesis} on the clause ``$m$ is generated from the $b$-particle''. On one manifold the
$b$-field is removable, tears the binding, or screens it without supplying a scale. Between manifolds it acts, but on frequencies, never on rates. Read in the dictionary of
\cref{s0:sec:question}, the rates belong to the symmetric part, the $c^2$ of the hypothesis. The directions of a Yang--Mills doublet that would
split them are not connections but symmetric terms. And the part of the copy field that survives every change of frame is the commutator
$[K,B]_{kl}=(\kappa_k-\kappa_l)B_{kl}$, proportional to the rate differences (\cref{d7:rem:verdict}). Read with $mc^2=\hbar\kappa$, so that a mass is a rate, the order of the hypothesis is
reversed in the arena, as it was for the Pati--Salam bosons of \cite{ES17}, which acquire mass rather than generate it. The symmetric part carries the rates.
Their differentiation, which is what a mass would be under $mc^2=\hbar\kappa$, breaks the copy symmetry. The $b$-field between the manifolds is the
gauge field of that broken symmetry, and it has an effect only because the symmetry is broken. Formally this is the structure of the
Englert--Brout--Higgs--Kibble mechanism with the rate matrix in the role of an adjoint vacuum value (interpretive: here $b$ has no kinetic term
and the rates are inputs).

What remains open is the origin of the rate differences. Within the arena there are two candidates: a rotor generator in $\uu n$ that is not
central, which is admitted but not selected, and a quantum scale for the $b$-field by dimensional transmutation (\cref{d4:conj:transmute}),
which is conjectural.

\begin{table}[t]
\centering
\small
\begin{tabular}{@{}>{\raggedright\arraybackslash}p{4.2cm}>{\raggedright\arraybackslash}p{6.6cm}>{\raggedright\arraybackslash}p{3.4cm}@{}}
\toprule
statement & what the arena gives & verdict\\
\midrule
the algebra is duplicated & $\Cl(6,6)=\Cl(3,3)\hat\otimes\Cl(3,3)$, Klein factor, Fock space (\cref{d5:sec:algebra}) & forced; inherited\\
the energy is duplicated & a beam splitter $S^3$, an $\eta$-circle; copying excluded (\cref{d5:thm:dup}) & chosen; cloning obstructed\\
the duplicate adds structure & complex structure, copy charge, $\so{2,\C}$, $S^4$ vacuum, balance surface & fixed (partly restated)\\
a pair from one null field & impossible covariantly; two null fields or a third body (\cref{d6:prop:covariant,d6:thm:pair}) & obstructed; inherited\\
the mass of one copy & not invariant under the copy rotation; gauge dependent if it is gauged (\cref{d6:prop:gauge}) & fixed; distinction open\\
a $b$-field between manifolds & $\so{n,\C}$, compact $\so n$; spectator manifold decoupled (\cref{d7:prop:exist}) & fixed\\
it is a deal breaker & non-removable at unequal rates; seesaw, confinement, resonance; time-odd at $n=3$ & fixed\\
it generates a rate & rates are gauge invariants (\cref{d7:thm:schr}(iii)) & obstructed\\
the rates differ & not in the paper's closure (common rate imposed); a non-central rotor in $\uu n$ (\cref{d7:prop:rotor}) & fixed / constructed; origin open\\
isospin one-half between manifolds & integral copy isospin at $n=2,3$; doublet $=$ one connection $+$ two rate terms; copy doublets from $n=4$ & obstructed at $n=2,3$ (common action)\\
the arrow of the hypothesis & $[K,B]\propto$ rate differences makes $b$ physical & reversed under $mc^2=\hbar\kappa$; interpretive\\
\bottomrule
\end{tabular}
\caption{Verdicts on the second question.}
\label{s6:tab:verdict}
\end{table}
